A diabetes mellitus (DM) tem previsão de aumento global substancial, passando de 463 milhões de indivíduos em 2019 para 700 milhões em 2045 \cite{teo_global_2021}. A retinopatia diabética (RD), complicação microvascular específica do DM \cite{wong_diabetic_2016, wang_diabetic_2018}, representa importante causa de cegueira evitável na população economicamente ativa e na terceira idade \cite{wong_diabetic_2016, wang_diabetic_2018, melo_current_2018, korn_malerbi_feasibility_2022, teo_global_2021}, com custos significativos para os indivíduos e os sistemas de saúde \cite{steinmetz_causes_2021}. Dados de 2020 estimam que 103,12 milhões de adultos convivam com RD, sendo 28,54 milhões com formas ameaçadoras à visão (VTDR). A tendência é de aumento: até 2045, esses números devem atingir 160,5 milhões e 44,82 milhões, respectivamente \cite{teo_global_2021}.

Este projeto objetiva desenvolver métodos avançados de IA para o auxílio diagnóstico dessas condições por meio da análise automatizada de imagens de retinografia, abordando tanto a identificação precoce das lesões quanto a avaliação de riscos metabólicos associados. Considerando-se que a RD apresenta fatores de risco controláveis, como hiperglicemia e hipertensão \cite{wong_diabetic_2016, melo_current_2018}, o diagnóstico precoce possibilita intervenções eficazes para impedir a progressão da doença.

A justificativa reside na capacidade da IA de reduzir custos, aumentar a acessibilidade ao diagnóstico precoce e melhorar a eficiência do atendimento no sistema público de saúde, especialmente em áreas remotas e com menor disponibilidade de especialistas. O uso de retinografia não-midriática, telemedicina e ferramentas automatizadas tem se mostrado promissor em diversos países \cite{wong_diabetic_2016, korn_malerbi_feasibility_2022, topol_high-performance_2019}, inclusive no Brasil, onde já foram testadas soluções com IA em contextos regionais como o estado de Sergipe \cite{korn_malerbi_feasibility_2022}.